\subsection{负范数估计}
\label{sec:ch05-negative-norm-estimates}
\setcounter{thmcount}{0}
\setcounter{equation}{0}

上一节作为副产物给出了 $L^2$ 误差估计, 它比 $H^1$ 误差估计多一个因子 $h$. 在某些情形下, 还能在更低的负范数中继续得到改进的误差估计. 负范数估计中出现更高次幂的 $h$, 可解释为误差具有振荡性. 回忆 $H^{-s}$ 范数定义为
\[
\MathBookFormulaID{formula-ch05-negative-sobolev-norm}
\|u\|_{H^{-s}(\Omega)}
=\sup_{0\ne v\in H^s(\Omega)}\frac{(u,v)}{\|v\|_{H^s(\Omega)}}.
\]
假设正则性估计
\MBSourceItem{1}
\begin{equation}
\MathBookFormulaID{formula-ch05-negative-norm-regularity}
\label{eq:ch05-negative-norm-regularity}
\|u\|_{H^{s+2}(\Omega)}\leq C_R\|f\|_{H^s(\Omega)}
\end{equation}
和逼近估计
\MBSourceItem{2}
\begin{equation}
\MathBookFormulaID{formula-ch05-negative-norm-approximation}
\label{eq:ch05-negative-norm-approximation}
\inf_{v\in V_h}\|u-v\|_{H^1(\Omega)}
\leq Ch^{s+1}\|u\|_{H^{s+2}(\Omega)}
\end{equation}
对某个 $s\geq0$ 成立.

\MBSourceItem{3}
\begin{Theorem}
\label{thm:ch05-negative-norm-error}
在条件~\eqref{eq:ch05-general-form-continuity}、\eqref{eq:ch05-shifted-coercivity}、\eqref{eq:ch05-negative-norm-regularity} 和~\eqref{eq:ch05-negative-norm-approximation} 下, 存在常数 $h_0,C$, 使对所有 $h\leq h_0$, 式~\eqref{eq:ch05-general-discrete-problem} 的解满足
\[
\MathBookFormulaID{formula-ch05-negative-norm-error}
\|u-u_h\|_{H^{-s}(\Omega)}
\leq Ch^{s+1}\|u-u_h\|_{H^1(\Omega)}.
\]
\end{Theorem}
\begin{Proof}
再次使用对偶性方法. 对任意 $\phi\in H^s(\Omega)$, 需要估计 $(u-u_h,\phi)$. 令 $w$ 为条件~\eqref{eq:ch05-negative-norm-regularity} 保证的伴随问题的解:
\MathBookSourcePage{159}
\[
\MathBookFormulaID{formula-ch05-negative-adjoint-problem}
a(v,w)=(v,\phi)\quad\forall v\in V.
\]
则对任意 $w_h\in V_h$,
\[
\MathBookFormulaID{formula-ch05-negative-duality-chain}
\begin{aligned}
(u-u_h,\phi)
&=a(u-u_h,w)\\
&=a(u-u_h,w-w_h)\qquad\text{(由式~\eqref{eq:ch05-galerkin-orthogonality})}\\
&\leq C_1\|u-u_h\|_{H^1(\Omega)}\|w-w_h\|_{H^1(\Omega)}.
\end{aligned}
\]
适当选择 $w_h$, 得
\[
\MathBookFormulaID{formula-ch05-negative-duality-regularity}
\begin{aligned}
(u-u_h,\phi)
&\leq Ch^{s+1}\|u-u_h\|_{H^1(\Omega)}\|w\|_{H^{s+2}(\Omega)}\\
&\leq Ch^{s+1}\|u-u_h\|_{H^1(\Omega)}\|\phi\|_{H^s(\Omega)}.
\end{aligned}
\]
对 $\phi$ 取上确界即得结论.
\end{Proof}

该定理说明逼近在正确值附近振荡, 因为形如 $(u-u_h,\phi)/\|\phi\|_{H^1(\Omega)}$ 的``加权平均''比 $u-u_h$ 本身多小一个因子 $h$.

负范数中改进收敛速度的关键要求是正则性估计~\eqref{eq:ch05-negative-norm-regularity}. 当它不成立时(见例~\ref{ex:ch05-nonconvex-sector-singularity}), 改进也不会发生. 用式~\eqref{eq:ch05-general-discrete-problem} 可作如下简单说明:
\[
\MathBookFormulaID{formula-ch05-pollution-energy-lower-bound}
(u-u_h,f)=a(u,u-u_h)=a(u-u_h,u-u_h)
\geq\alpha\|u-u_h\|_{H^1(\Omega)}^2.
\]
若 $f\in H^s(\Omega)$, 则
\MBSourceItem{4}
\begin{equation}
\MathBookFormulaID{formula-ch05-pollution-negative-lower-bound}
\label{eq:ch05-pollution-negative-lower-bound}
\begin{aligned}
\|u-u_h\|_{H^{-s}(\Omega)}
&=\sup_{0\ne\phi\in H^s(\Omega)}
\frac{(u-u_h,\phi)}{\|\phi\|_{H^s(\Omega)}}\\
&\geq\frac{(u-u_h,f)}{\|f\|_{H^s(\Omega)}}
\geq\alpha\frac{\|u-u_h\|_{H^1(\Omega)}^2}{\|f\|_{H^s(\Omega)}}.
\end{aligned}
\end{equation}
因此误差的负范数绝不可能小于其 $H^1(\Omega)$ 范数平方的量级. 这通常称为边界奇点的``污染效应''(参见 Ciarlet \& Lions 1991 中 Wahlbin 的文章).

下面证明对例~\ref{ex:ch05-nonconvex-sector-singularity} 中的问题, 有 $\|u-u_h\|_{H^1(\Omega)}\geq ch^\beta$. 一般而言, 若 $V_h$ 由次数不超过 $k+1$ 的分片多项式组成, 则
\MathBookSourcePage{160}
\[
\MathBookFormulaID{formula-ch05-local-polynomial-error-lower-bound}
\begin{aligned}
|u-u_h|_{H^1(\Omega)}^2
&\geq\inf_{v\in V_h}\sum_{T\in\mathcal T^h}\int_T|\nabla u-\nabla v|^2\,dx\\
&\geq\sum_{T\in\mathcal T^h}\inf_{\widetilde v\in P_k^2}
\int_T|\nabla u-\widetilde v|^2\,dx\\
&\geq\inf_{\widetilde v\in P_k^2}\int_{\widetilde T}|\nabla u-\widetilde v|^2\,dx
\end{aligned}
\]
对任意 $\widetilde T\in\mathcal T^h$ 成立. 这里 $P_k$ 表示次数不超过 $k$ 的多项式. 取一个以原点为顶点的单元 $\widetilde T$, 由齐次性论证(练习~\ref{exr:ch05-21} 和~\ref{exr:ch05-22})得
\[
\MathBookFormulaID{formula-ch05-sector-homogeneous-lower-bound}
\inf_{\widetilde v\in P_k^2}\int_{\widetilde T}
|\nabla(r^\beta\sin\beta\theta)-\widetilde v|^2\,dx
\geq c_0\operatorname{diam}(\widetilde T)^{2\beta}.
\]
因为 $u-r^\beta\sin\beta\theta=-r^{2+\beta}\sin\beta\theta$, 由三角不等式有
\[
\MathBookFormulaID{formula-ch05-sector-solution-lower-bound}
\inf_{\widetilde v\in P_k^2}\int_{\widetilde T}|\nabla u-\widetilde v|^2\,dx
\geq c_0\operatorname{diam}(\widetilde T)^{2\beta}
-c_1\operatorname{diam}(\widetilde T)^{2+2\beta}.
\]
若 $\mathcal T^h$ 拟一致, 则
\MBSourceItem{5}
\begin{equation}
\MathBookFormulaID{formula-ch05-sector-global-error-lower-bound}
\label{eq:ch05-sector-global-error-lower-bound}
\|u-u_h\|_{H^1(\Omega)}\geq c_2h^\beta,
\end{equation}
其中 $c_2>0$, 只要 $h$ 足够小.
